\documentclass[12pt]{article}
\usepackage[utf8]{inputenc}
\usepackage[T1]{fontenc}
\usepackage[margin=1in]{geometry}
\usepackage{setspace}
\usepackage{xcolor}
\usepackage[colorlinks=true, urlcolor=blue!60!black, linkcolor=blue!60!black]{hyperref}
\setlength{\parindent}{2em}
\setlength{\parskip}{1.25em}
\renewcommand{\baselinestretch}{1.0}
\title{\textit{Craig v. Missouri}, 29 U.S. (4 Pet.) 410 (1830)}
\author{Case Brief}
\date{}
\begin{document}
\maketitle

\section{Case Metadata}

\begin{itemize}
  \item \textbf{Case Name:} \textit{Craig v. Missouri} (full caption: \textit{Hiram Craig, John Moore, and Ephraim Moore v. State of Missouri})
  \item \textbf{Citation:} \href{https://supreme.justia.com/cases/federal/us/29/410/}{29 U.S. (4 Pet.) 410 (1830)}
  \item \textbf{Court and Dates:} Supreme Court of the United States; argued March 2--3, 1830; decided March 12, 1830
  \item \textbf{Chapter:} Chapter 4: Federalism and State Sovereignty
  \item \textbf{Vote:} 4--3 (Majority: Marshall, C.J., joined by Duvall, Story, and Baldwin, JJ.; Dissenting: Johnson, Thompson, and McLean, JJ.)
  \item \textbf{Author of the Opinion of the Court:} Chief Justice John Marshall
  \item \textbf{Dissenting Opinions:} Justice William Johnson; Justice Smith Thompson; Justice John McLean
  \item \textbf{Constitutional and Statutory Provisions:} U.S. Const. art. I, \S~10, cl. 1 (Bills of Credit Clause); Judiciary Act of 1789, \S~25
  \item \textbf{Primary Source Links:} \href{https://supreme.justia.com/cases/federal/us/29/410/}{Justia U.S. Supreme Court Center}; \href{https://www.law.cornell.edu/constitution-conan/article-1/section-1/clause-1/coining-money-by-states}{Cornell LII Constitution Annotated}
\end{itemize}

\section{Facts and Procedural History}

In the wake of the Panic of 1819 and a severe shortage of specie on the western frontier, the newly admitted State of Missouri enacted ``An act for the establishment of loan offices'' on June 27, 1821. The statute authorized the state auditor and treasurer to issue up to \$200,000 in paper ``certificates'' in small denominations ranging from fifty cents to ten dollars, bearing 2\% annual interest. Each certificate stated on its face that it ``shall be receivable at the Treasury, or any of the loan offices of the State of Missouri, in the discharge of taxes or debts due to the state,'' and the act further made the certificates payable for salaries and fees of state officers and receivable by lessees of the state's public salt springs in payment for salt.

Loan-office commissioners were empowered to lend these certificates to Missouri citizens on real-estate mortgages or, for sums under \$200, on personal security. Section 23 pledged a dedicated redemption fund---consisting of proceeds from the public salt springs, accruing loan interest, estates acquired under the act, and debts owed to the state---and expressly declared that ``the faith of the state is hereby also pledged for the same purpose.'' One-tenth of the certificates were to be withdrawn annually. Although the act did not declare the certificates legal tender for private debts, they were designed to circulate as a local medium of exchange.

On August 1, 1822, Hiram Craig, John Moore, and Ephraim Moore executed a joint promissory note to the State of Missouri at the Chariton loan office for \$199.99, plus 2\% annual interest, in exchange for loan-office certificates. After they defaulted, Missouri sued them in assumpsit in state circuit court. The defendants pleaded \textit{non assumpsit}, arguing that the certificates were unconstitutional ``bills of credit'' under Article I, Section 10, Clause 1, rendering the note's consideration illegal and void. The circuit court entered judgment for Missouri, and the Supreme Court of Missouri affirmed. Craig and the Moores then brought a writ of error to the U.S. Supreme Court under Section 25 of the Judiciary Act of 1789, where Senator Thomas Hart Benton argued the case for Missouri.

\section{Issue(s) Presented}

\begin{enumerate}
  \item \textbf{Appellate Jurisdiction:} Did the Supreme Court have jurisdiction under Section 25 of the Judiciary Act of 1789 when the Supreme Court of Missouri's judgment did not explicitly state on the face of the record that it was upholding the state statute against a federal constitutional challenge?
  \item \textbf{Bills of Credit Clause:} Did Missouri's interest-bearing loan-office certificates---issued directly by state officers on the pledged faith of the state to circulate as money, but not made legal tender for private debts---constitute prohibited ``bills of credit'' under Article I, Section 10, Clause 1?
  \item \textbf{Enforceability of the Note:} If the certificates were unconstitutional bills of credit, was Craig's promissory note given in exchange for them void for illegal consideration?
\end{enumerate}

\section{Holding}

\textbf{Reversed (4--3).} Chief Justice Marshall held for the Court that: (1) Section 25 jurisdiction was proper because the record showed that the federal constitutional validity of Missouri's loan-office statute was necessarily decided below; (2) Missouri's loan-office certificates were unconstitutional ``bills of credit'' under \href{https://www.law.cornell.edu/constitution-conan/article-1/section-1/clause-1/coining-money-by-states}{Article I, Section 10, Clause 1} because they were paper issued by the state, on the faith of the state, intended to circulate through the community as money---regardless of their statutory label or lack of private legal-tender status; and (3) because the certificates were emitted in violation of the Constitution, the promissory note given for them rested on illegal consideration and was void.

\section{Reasoning}

\textbf{1. Section 25 Jurisdiction.} Rejecting Missouri's threshold objection, Marshall explained that a state court cannot insulate federal questions from Supreme Court review merely by omitting an explicit recital that a challenged state statute is constitutional. Relying on \textit{Martin v. Hunter's Lessee} (1816) and \textit{Wilson v. Black Bird Creek Marsh Co.} (1829), the Court inspected the special statement of facts in the record and found that the sole defense turned on the constitutionality of the loan-office act; affirming judgment for Missouri necessarily upheld the statute against the federal Constitution.

\textbf{2. Definition of ``Bills of Credit.''} Turning to the merits, Marshall construed ``emit bills of credit'' in light of pre-1787 colonial and Revolutionary usage. The phrase did not bar ordinary state borrowing or contracts to pay for goods and services; rather, it prohibited a state government from issuing a paper medium intended to circulate through the community for its ordinary purposes as money, redeemable at a future day. Missouri's certificates exhibited every hallmark of such paper currency: they were issued directly by state officials in small, spendable denominations (\$0.50 to \$10), lent out to enter circulation, made receivable for taxes, public debts, salaries, and salt, and expressly backed by the pledged funds and ``faith of the state.''

\textbf{3. Substance Over Form.} Missouri argued that its statute created ``loan-office certificates,'' not ``bills of credit.'' Marshall rejected this formalist evasion outright, reasoning that the Constitution prohibits things, not mere labels. Giving a new name to an old mischief could not rescue a state paper currency from Article I, Section 10.

\textbf{4. Independence from the Legal Tender Clause.} Missouri also contended that paper could not be a ``bill of credit'' unless made legal tender for private debts. Marshall dismantled this argument on both textual and historical grounds. Textually, Article I, Section 10, Clause 1 contains two separate prohibitions: no state shall ``emit Bills of Credit,'' \textit{or} ``make any Thing but gold and silver Coin a Tender in Payment of Debts.'' Reading legal-tender status into ``bills of credit'' would render the Bills of Credit Clause redundant and effectively erase it from the Constitution. Historically, Marshall surveyed colonial emissions---including Virginia's treasury notes of 1769, 1771, 1773, and 1775---which were not legal tender yet circulated alongside tender bills as ``bills of credit'' and produced the same inflationary evils.

\textbf{5. Void Consideration.} Because Missouri's emission of the certificates violated a direct constitutional prohibition, the transaction was unlawful. A court cannot enforce a promissory note whose sole consideration is an unconstitutional act; accordingly, judgment had to be entered for Craig and the Moores.

\section{Concurring and Dissenting Opinions}

No separate concurring opinion was published in 1830. However, seven years later in \href{https://supreme.justia.com/cases/federal/us/36/257/}{\textit{Briscoe v. Bank of Commonwealth of Kentucky}, 36 U.S. (11 Pet.) 257 (1837)}, Justice Henry Baldwin revealed that he had provided the decisive fourth vote in \textit{Craig} on narrower grounds than Marshall's opinion: Baldwin had requested the sentence noting that Missouri pledged ``the faith and funds of the state'' and believed a bill of credit required both a pledge of state faith and tender status for some debts.

All three dissenters in \textit{Craig} wrote separately, urging narrower definitions of ``bills of credit'' that would preserve state fiscal flexibility:
\begin{itemize}
  \item \textbf{Justice William Johnson (dissenting):} Offered two functional distinctions. First, because Missouri's certificates bore 2\% interest, their value changed every day, disqualifying them from acting as a uniform circulating medium. Second, whereas Revolutionary bills were unsecured promises to \textit{pay}, Missouri's certificates were promises to \textit{receive} in payment of taxes and state debts, effectively drawing on an enforceable revenue fund rather than bare sovereign faith.
  \item \textbf{Justice Smith Thompson (dissenting):} Defined a constitutional bill of credit as ``a bill drawn and resting merely on the credit of the drawer, as contradistinguished from a fund constituted or pledged for the payment of the bill'' (quoted in \href{https://supreme.justia.com/cases/federal/us/36/257/}{\textit{Briscoe}, 36 U.S. at 313--314}). Because Missouri appropriated specific funds (salt springs, interest, and state debts) and did not force private parties to accept the certificates, the act was constitutional.
  \item \textbf{Justice John McLean (dissenting):} Argued that a bill of credit must be issued by a state with its circulation enforced by statutory compulsion and without an appropriated redemption fund. Emphasizing judicial restraint and federalism, McLean cautioned that federal and state powers intermix ``like the colours of the rainbow'' and require ``mutual forbearance'' before invalidating state legislation.
\end{itemize}

\section{Key Quotations}

\begin{quote}
``To emit bills of credit conveys to the mind the idea of issuing paper, intended to circulate through the community for its ordinary purposes as money, which paper is redeemable at a future day; this is the sense in which it has been always understood.''
--- Chief Justice Marshall, \href{https://supreme.justia.com/cases/federal/us/29/410/}{29 U.S. at 432} (quoted verbatim in \href{https://supreme.justia.com/cases/federal/us/36/257/}{\textit{Briscoe}, 36 U.S. at 329}).
\end{quote}

\begin{quote}
``Is the proposition to be maintained that the Constitution meant to prohibit names, and not things? That a very important act, big with great and ruinous mischief, which is expressly forbidden by words most appropriate for its description may be performed by the substitution of a name? That the Constitution, in one of its most important provisions, may be openly evaded by giving a new name to an old thing? We cannot think so.''
--- Chief Justice Marshall, \href{https://supreme.justia.com/cases/federal/us/29/410/}{29 U.S. at 433} (quoted verbatim in \href{https://supreme.justia.com/cases/federal/us/36/257/}{\textit{Briscoe}, 36 U.S. at 332}).
\end{quote}

\begin{quote}
``The Constitution therefore considers the emission of bills of credit and the enactment of tender laws as distinct operations, independent of each other, which may be separately performed; both are forbidden. To sustain the one because it is not also the other---to say, that bills of credit may be emitted if they be not a tender of debts---is in effect to expunge that distinct, independent prohibition and to read the clause as if it had been entirely omitted. We are not at liberty to do this.''
--- Chief Justice Marshall, \href{https://supreme.justia.com/cases/federal/us/29/410/}{29 U.S. at 434} (quoted verbatim in \href{https://supreme.justia.com/cases/federal/us/36/257/}{\textit{Briscoe}, 36 U.S. at 333}).
\end{quote}

\begin{quote}
``To constitute a `bill of credit' within the meaning of the Constitution, it must be issued by a state, and its circulation as money enforced by statutory provisions. It must contain a promise of payment by the state generally when no fund has been appropriated to enable the holder to convert it into money.''
--- Justice McLean (dissenting), \href{https://supreme.justia.com/cases/federal/us/29/410/}{29 U.S. at 453} (quoted verbatim in \href{https://supreme.justia.com/cases/federal/us/36/257/}{\textit{Briscoe}, 36 U.S. at 314}).
\end{quote}

\section{Doctrinal Significance and Modern Legacy (Through 2026)}

\textit{Craig v. Missouri} marked the high-water mark of the Marshall Court's nationalist enforcement of Article I, Section 10's monetary restraints, yet its practical reach was swiftly curtailed by the Taney Court. After Marshall's death and Andrew Jackson's judicial appointments reshaped the bench, the Court decided \href{https://supreme.justia.com/cases/federal/us/36/257/}{\textit{Briscoe v. Bank of Commonwealth of Kentucky}, 36 U.S. (11 Pet.) 257 (1837)} (6--1). Writing for the majority, \textit{Craig} dissenter Justice McLean upheld circulating bank notes issued by a wholly state-owned, legislature-controlled bank. Purporting to leave \textit{Craig} intact, \textit{Briscoe} drew a sharp formal distinction: because the bank was a separately incorporated, suable entity issuing notes on its own corporate funds rather than directly on the faith of the state, its notes were not ``bills of credit.'' Justice Story, the sole surviving member of the \textit{Craig} majority who dissented, protested that \textit{Briscoe} allowed states to do indirectly what \textit{Craig} barred directly, leaving the constitutional prohibition a hollow shell.

Later nineteenth-century decisions further narrowed \textit{Craig}. In \href{https://www.law.cornell.edu/supremecourt/text/114/270}{\textit{Poindexter v. Greenhow}, 114 U.S. 270 (1885)} (the Virginia Coupon Cases), the Court held that tax-receivable interest coupons detached from state bonds were not bills of credit because they were not designed to circulate as money. Similarly, in \href{https://www.law.cornell.edu/supremecourt/text/177/66}{\textit{Houston \& Texas Central R.R. Co. v. Texas}, 177 U.S. 66 (1900)}, the Court upheld state treasury warrants used to pay public debts, demanding clear proof of legislative intent that the paper circulate as currency.

Although the Civil War National Bank Acts, the Federal Reserve System, and federal legal-tender supremacy rendered state paper money largely obsolete for over a century, \textit{Craig}'s framework has regained unexpected salience through 2026 in two arenas. First, states have increasingly experimented with precious-metal tender statutes and state-sponsored digital payment instruments, prompting renewed scholarly scrutiny of Article I, Section 10's \href{https://constitution.heritage.org/essays/a1-s10-c1-c/}{State Money Clause}. Second, on January 7, 2026, the \href{https://stabletoken.wyo.gov/}{Wyoming Stable Token Commission} launched the Frontier Stable Token (FRNT) for public purchase---the nation's first state-issued, fiat-backed stablecoin. Deliberately echoing \textit{Briscoe}'s safe harbor, Wyoming structured the Commission as a separate ``body politic and corporate'' and mandated \href{https://stabletoken.wyo.gov/faqs}{102\% statutory reserves} in segregated U.S. dollars and short-term Treasuries rather than an unbacked pledge of state faith. Whether modern state-issued digital currencies fall within \textit{Craig}'s anti-evasion ban or \textit{Briscoe}'s corporate-reserve exception remains a live constitutional question in 2026.

\section{Pedagogical Critiques and Exam Hypotheticals}

\textbf{Pedagogical Critiques.}
\begin{itemize}
  \item \textbf{Formalism Versus Functionalism:} Placed alongside \textit{Briscoe}, \textit{Craig} illustrates a classic constitutional law tension. Marshall insisted that the Constitution bars ``things,'' not ``names,'' yet \textit{Briscoe} permitted states to achieve virtually identical circulating currency by interposing a state-owned corporation---inviting students to ask whether structural limits in Article I, Section 10 survive corporate form.
  \item \textbf{Monetary History and Frontier Realism:} Economic historians and the \textit{Craig} dissenters note that Missouri's capped (\$200,000), interest-bearing, non-tender certificates were a modest response to acute specie scarcity rather than the reckless, forced-tender paper money of 1786 Rhode Island. Marshall's categorical rule prioritized uniform national monetary discipline over state-level macroeconomic experimentation.
  \item \textbf{Litigant Windfall:} Because the remedy was invalidating the underlying contract for illegal consideration, Craig and the Moores received value from the loan-office certificates and escaped repayment altogether---a recurring dilemma when structural constitutional limits are enforced defensively in private debt litigation.
\end{itemize}

\textbf{Exam Hypothetical 1 (Direct \textit{Craig} Violation).} Facing a severe budget impasse, the State of Ozark enacts a statute directing its State Treasurer to print \$500 million in ``Fiscal Liquidity Certificates'' in \$5, \$20, and \$100 bearer denominations. The certificates are not legal tender for private debts, but they are receivable at face value for all state taxes and public-university tuition, and the statute pledges ``the full faith and credit of the State of Ozark'' for redemption in five years. \textit{Analysis:} Unconstitutional under \textit{Craig}. Even without private legal-tender status or the label ``bills of credit,'' the certificates are issued directly by state officers, backed by the state's faith, and structured in small bearer denominations to circulate as a common medium of exchange.

\textbf{Exam Hypothetical 2 (``Opposite Side'' Contrast under \textit{Briscoe} and 2026 Digital Assets).} By contrast, suppose the State of Frontier charters an independent, suable public corporation to issue ``Frontier Digital Dollars''---blockchain tokens redeemable on demand for \$1.00 U.S., backed 102\% by segregated U.S. dollars and short-term Treasury bills held in trust, with an explicit statutory disclaimer that neither the state's general faith nor its taxing power backs the tokens. \textit{Analysis:} Under \textit{Briscoe} and \textit{Houston \& Texas Central}, the tokens likely survive Article I, Section 10 scrutiny because they are issued by a separately suable corporate instrumentality, backed by a dedicated, fully liquid reserve fund rather than sovereign credit, and redeemable on demand. Conversely, a functionalist challenger invoking Marshall's ``names, and not things'' principle in \textit{Craig} and Story's \textit{Briscoe} dissent would argue that a state-created, state-governed entity minting a circulating digital medium designed for everyday commerce and taxes is an impermissible indirect emission of bills of credit.

\end{document}
